\chapter{多元函数微分}
\begin{definition}{向量和多元函数}{}
	$x_1,x_2,\ldots,x_n \in \RR$向量
	\[
		X := [x_1,x_2,\ldots,x_n]\transp,
	\]
	定义在开集$\Omega \subset \RR^n$上的$n$元函数$f:\Omega \to \RR.$
\end{definition}
$f_1,f_2,\ldots,f_m:\Omega \to \RR$构成向量值函数$F = [f_1,f_2,\ldots,f_m]\transp.$
\section{偏导数}
\begin{definition}{}{}
	$\uvec{e}_1,\uvec{e}_2,\ldots,\uvec{e}_n$ 是$\RR^n$的标准正交基，$X_0 \in \Omega$，则$f$在$X_0$点关于$x_i$的偏导数
	\begin{equation}\label{eq:calc-partial-derivative}
		\pv f{x_i}(X_0) := \lim_{t\to 0}\frac{f(X_0 + t\uvec{e}_i) - f(X_0)}t,
	\end{equation}
	若$f$在$X_0$可微，其全微分为
	\begin{align}
		\d f(X_0) = \sum_{i=1}^n\pv f{x_i}(X_0)\d x_i.
	\end{align}
\end{definition}
有时用$\p_if$比$\p f/\p x_i$更好.
\begin{theorem}{可微性判定}{}
	对二元函数$f(x,y)$
	\begin{center}
		两个一阶偏导在$X_0$的邻域内存在，且在$X_0$连续
		\vimplies
		$f$在$X_0$可微
		\viff
		$f(X) - \Bigb{f(X_0) + \pv fx(X_0)\Delta x + \pv fy(X_0)\Delta y} = o(\norm{\Delta X})$
	\end{center}
	其中$\Delta X \equiv [\Delta x,\Delta y]\transp.$
\end{theorem}
\section{方向导数和梯度}
\begin{definition}{}{}
	$f$可微，其沿单位向量$\uvec{\ell} = [\cos\alpha_1,\cos\alpha_2,\ldots,\cos\alpha_n]\transp$方向的方向导数 % \bm \ell=[x_1,x_2,\ldots,x_n], \bm \ell/|\ell|=
	\begin{equation}\label{eq:calc-directional-derivative}
		\pv f{\uvec{\ell}}(X_0) := \lim_{t\to 0}\frac{f(X_0 + t\uvec{\ell}) - f(X_0)}t;
	\end{equation}
	梯度
	\begin{equation}\label{eq:calc-gradient}
		\grad f := \Bigb{\pv f{x_1},\pv f{x_2},\ldots,\pv f{x_n}}\transp.
	\end{equation}
\end{definition}
因此方向导数也可以看做
\begin{align}
	\pv f{\uvec{\ell}}(X_0) = \sum_{i=1}^n\pv f{x_i}(X_0)\cos\alpha_i = \grad f(X_0) \cdot \uvec{\ell}.
\end{align}
\section{高阶偏导数}
\begin{definition}{}{}
	若$\pv f{x_i}$对$x_j$的偏导存在，则记二阶偏导数
	\begin{subequations}\label{eq:calc-second-partials}
		\begin{align}
			\pw f{x_j}{x_i} &:= \pp{x_j}\pv f{x_i},\\
			\pv[2]f{x_i} &\equiv \pw f{x_i}{x_i}.
		\end{align}
	\end{subequations}
	第一行的二阶偏导数也可记作$\p_{ji}f$。
\end{definition}
\begin{theorem}{混合偏导交换的充分条件}{}
	\begin{center}
		$f$在$X_0$的邻域内二阶连续可微
		\vimplies
		$\pw f{x_i}{x_j} = \pw f{x_j}{x_i}$
	\end{center}
\end{theorem}
\section{向量值函数微分}
\begin{definition}[breakable=false]{Jacobi矩阵}{}
	可微的向量值函数$F:\Omega \subset \RR^n \to \RR^m$%，于是$\bm f:\Omega\to\RR^m$
	\begin{equation}\label{eq:calc-vector-differential}
		\begin{gathered}
			\begin{bmatrix}
				\d f_1 \\
				\vdots     \\
				\d f_m
			\end{bmatrix} =
			%\begin{bmatrix}
			%    \frac{\p f_1}{\p x_1}\d x_1+\cdots+\frac{\p f_1}{\p x_n}\d x_n \\
			%    \vdots                                                         \\
			%    \frac{\p f_m}{\p x_1}\d x_1+\cdots+\frac{\p f_m}{\p x_n}\d x_n
			%\end{bmatrix}=
			\begin{bmatrix}
				\pv{f_1}{x_1} & \cdots & \pv{f_1}{x_n} \\
				\vdots                & \ddots & \vdots                \\
				\pv{f_m}{x_1} & \cdots & \pv{f_m}{x_n}
			\end{bmatrix}
			\begin{bmatrix}
				\d x_1 \\
				\vdots \\
				\d x_n
			\end{bmatrix};\\
			\d F = J_F\d X.
		\end{gathered}
	\end{equation}
	式中的$m \times n$矩阵称为Jacobi矩阵
	\begin{equation}\label{eq:calc-jacobian-matrix}
		J_F \equiv \frac{\p(f_1,\ldots,f_m)}{\p(x_1,\ldots,x_n)}.
	\end{equation}
	$m = n$时，对应有Jacobi行列式
	\begin{equation}\label{eq:calc-jacobian-determinant}
		\frac{D(f_1,\ldots,f_n)}{D(x_1,\ldots,x_n)} := \det J_F
	\end{equation}
\end{definition}
\section{复合函数微分}
\begin{theorem}{}{}
	$G:\RR^n \to \RR^m,F:\RR^m \to \RR^k$，$Y_0 = G(X_0)$。若$G$在$X_0$可微、$F$在$Y_0$可微，则复合函数在$X_0$可微，链式法则有以下等价表达
	\begin{compactenum}[(1)]
		\item $\d(F \circ G)(X_0) = \d F(Y_0) \circ \d G(X_0);$
		\item $J_{F\circ G}(X_0) = J_F(Y_0)J_G(X_0);$
		\item 用各分量的偏导数表示，
		\begin{equation}\label{eq:calc-jacobian-chain-rule}
			\frac{\p(f_1,\ldots,f_k)}{\p(x_1,\ldots,x_n)}(X_0)
			= \frac{\p(f_1,\ldots,f_k)}{\p(y_1,\ldots,y_m)}(Y_0)
			 \frac{\p(g_1,\ldots,g_m)}{\p(x_1,\ldots,x_n)}(X_0).
		\end{equation}
	\end{compactenum}
\end{theorem}
特别地，当$k = 1$即$F = f$为标量值函数时
\begin{align}
	\pv{f \circ G}{x_i}(X_0) = \sum_{j=1}^m\pv f{y_j}(Y_0)\pv{g_j}{x_i}(X_0).
\end{align}
$\color{gray}\text{例：}\pv{f(x,y,x^2)}x = \p_1f(x,y,x^2) + 2x\,\p_3f(x,y,x^2)$
\section{隐函数}
本节采用常见的$\cont^1$充分条件，参见\href{https://math.rice.edu/~hardt/401F03/ift.pdf}{Rice University：Inverse and Implicit Function Theorems}。
\begin{theorem}{二元隐函数}{}
    设$f$在$P = (x_0,y_0)$的邻域内属于$\cont^1$（一阶连续可微），且$f(P) = 0$、$f_y(P) \neq 0$。则在该点附近，方程$f(x,y) = 0$唯一确定经过$P$的局部$\cont^1$函数$y = y(x)$，且
    \begin{equation}\label{eq:calc-implicit-scalar-derivative}
        y'(x) = -\frac{f_x(x,y(x))}{f_y(x,y(x))}.
    \end{equation}
\end{theorem}
隐函数的存在性和可微性由定理保证；在此基础上，对恒等式$f(x,y(x)) = 0$使用链式法则，才得到上述求导公式。这不是存在性定理本身的证明。
\begin{theorem}{多元隐函数}{}
    设$f(X,y)$在$P = (X_0,y_0)$的邻域内属于$\cont^1$，其中$X \in \RR^n$，且$f(P) = 0$、$f_y(P) \neq 0$。则存在唯一的局部$\cont^1$函数$y = y(X)$，满足$y(X_0) = y_0$和$f(X,y(X)) = 0$，并且
    \begin{equation}\label{eq:calc-implicit-gradient}
        \pv y{x_i}(X) = -\frac{f_{x_i}(X,y(X))}{f_y(X,y(X))}.
    \end{equation}
\end{theorem}
\begin{theorem}{向量值隐函数}{}
    设$F = (f_1,\ldots,f_m)\transp$在$P = (X_0,Y_0) \in \RR^{n+m}$的邻域内属于$\cont^1$，$F(P) = \bm0$，且$m \times m$矩阵$J_YF(P)$可逆。则存在唯一的局部$\cont^1$函数$Y = Y(X)$，满足$Y(X_0) = Y_0$和$F(X,Y(X)) = \bm0$，并且
    \begin{equation}\label{eq:calc-implicit-jacobian}
        J_Y(X) = -\bigb{J_YF(X,Y(X))}\inv J_XF(X,Y(X)).
    \end{equation}
    这里$J_Y(X)$表示映射$X \mapsto Y(X)$的$m \times n$Jacobi矩阵，$J_XF,J_YF$分别表示$F$对两组自变量的偏导矩阵。
\end{theorem}
\begin{theorem}{反函数}{}
    设$F:\Omega \subset \RR^n \to \RR^n$属于$\cont^1$，$X_0 \in \Omega$且$\det J_F(X_0) \neq 0$。则$F$在$X_0$附近有$\cont^1$局部反函数。令$Y = F(X)$，在相应邻域内有
    \begin{equation}\label{eq:calc-inverse-jacobian}
        J_{F\inv}(Y) = \bigb{J_F(F\inv(Y))}\inv = \bigb{J_F(X)}\inv.
    \end{equation}
    Jacobi行列式非零保证的是局部可逆，不自动保证全局一一对应。
\end{theorem}
\begin{example}{二阶偏导数举例}{}
    设$f \in \cont^2$，$f_z \neq 0$，由$f(x,y,z) = 0$局部确定$z = z(x,y)$。由式\eqref{eq:calc-implicit-gradient}得到
    \[
        z_x = -\frac{f_x}{f_z},\qquad z_y = -\frac{f_y}{f_z}.
    \]
    再对恒等式$f_x + f_z z_x = 0$关于$y$求导，须包含$z(x,y)$的链式贡献：
    \[
        f_{xy} + f_{xz}z_y + (f_{zy} + f_{zz}z_y)z_x + f_z z_{xy} = 0.
    \]
    因而
    \begin{equation}\label{eq:calc-implicit-mixed-derivative}
        z_{xy} = -\frac{f_z^2f_{xy} - f_yf_zf_{xz} - f_xf_zf_{yz} + f_xf_yf_{zz}}{f_z^3}.
    \end{equation}
    右侧的$f$及其偏导均在$(x,y,z(x,y))$处取值。
\end{example}
\section{法与切}
给定非零向量$\bm v = [a,b,c]\transp$与点$P_0(x_0,y_0,z_0)$，可以确定
\begin{subequations}\label{eq:calc-line-plane-cartesian}
	\begin{align}
		\text{线}:   & \qquad \frac{x - x_0}a = \frac{y - y_0}b = \frac{z - z_0}c \\
		\text{平面}: & \qquad a(x - x_0) + b(y - y_0) + c(z - z_0) = 0
	\end{align}
\end{subequations}
上述直线的对称式要求$a,b,c$均非零；分量为零时改用参数式。平面的两条参数方向必须线性无关：
\begin{subequations}\label{eq:calc-line-plane-parametric}
	\begin{gather}
		\text{线}:~
		\begin{cases}
			x = x_0 + at \\[-1ex]
			y = y_0 + bt \\[-1ex]
			z = z_0 + ct
		\end{cases}
		\quad
		\text{平面}:~
		\begin{cases}
			x = x_0 + a_1u + b_1v \\[-1ex]
			y = y_0 + a_2u + b_2v \\[-1ex]
			z = z_0 + a_3u + b_3v
		\end{cases}
		\\
		a =
		\begin{vmatrix}
			a_2 & b_2 \\
			a_3 & b_3
		\end{vmatrix}\quad b =
		\begin{vmatrix}
			a_3 & b_3 \\
			a_1 & b_1
		\end{vmatrix}\quad c =
		\begin{vmatrix}
			a_1 & b_1 \\
			a_2 & b_2
		\end{vmatrix}.
	\end{gather}
\end{subequations}
\paragraph{曲面的法线和切平面}
以下函数至少属于$\cont^1$；参数曲面要求两个切向量线性无关，隐式曲面要求$\grad f(P_0) \neq \bm0$。
\subparagraph{显函数}
\begin{equation}\label{eq:calc-graph-normal}
	z = f(x,y) \implies \bm v = \Bigb{\pv fx,\pv fy,-1}_{(x_0,y_0)}.
\end{equation}
\subparagraph{参数方程}
\begin{equation}\label{eq:calc-parametrized-surface-normal}
	\begin{cases}
		x = f_1(u,v) \\
		y = f_2(u,v) \\
		z = f_3(u,v)
	\end{cases}
	\implies \bm v = \biggb{\frac{D(f_2,f_3)}{D(u,v)},\frac{D(f_3,f_1)}{D(u,v)},\frac{D(f_1,f_2)}{D(u,v)}}_{(u_0,v_0)}.
\end{equation}
\subparagraph{隐函数}
\begin{equation}\label{eq:calc-level-surface-normal}
	f(x,y,z) = 0 \implies \bm v = \Bigb{\pv fx,\pv fy,\pv fz}_{P_0} = \grad f(P_0).
\end{equation}
\paragraph{曲线的法平面和切线}
参数曲线要求$\bm r'(t_0) \neq \bm0$；两曲面交线要求$\grad F_1(P_0),\grad F_2(P_0)$线性无关。以下向量为曲线的切向量，也是法平面的法向量。
\subparagraph{参数}
\begin{equation}\label{eq:calc-space-curve-tangent}
	\begin{cases}
		x = x(t) \\[-1ex]
		y = y(t) \\[-1ex]
		z = z(t)
	\end{cases}
	\implies \bm v = [x'(t),y'(t),z'(t)].
\end{equation}
\subparagraph{隐函数}
\begin{equation}\label{eq:calc-intersection-tangent}
	\begin{cases}
		F_1(x,y,z) = 0 \\
		F_2(x,y,z) = 0
	\end{cases}
	\implies \bm v = \biggb{\frac{D(F_1,F_2)}{D(y,z)},\frac{D(F_1,F_2)}{D(z,x)},\frac{D(F_1,F_2)}{D(x,y)}}_{P_0}
\end{equation}
\section{Taylor公式}
\begin{definition}{Hesse矩阵}{}
	$f:\Omega \to \RR$的Jacobi矩阵是行向量$J_f = (\p_1f,\ldots,\p_nf)$。
	二阶偏导组成Hesse矩阵
	\begin{equation}\label{eq:calc-hessian}
		H_f :=
		\begin{bmatrix}
			\frac{\p^2f}{\p x_1^2}     & \cdots & \frac{\p^2f}{\p x_1\p x_n} \\
			\vdots                     & \ddots & \vdots                     \\
			\frac{\p^2f}{\p x_n\p x_1} & \cdots & \frac{\p^2f}{\p x_n^2}
		\end{bmatrix}.
	\end{equation}
\end{definition}
\paragraph{带Lagrange余项的1阶Taylor公式}
设$f \in \cont^2$于包含线段$[X_0,X]$的开邻域，$\Delta X = X - X_0$。则$\exists\theta \in (0,1)$，记$X_\theta := X_0 + \theta\Delta X$，有
\begin{align}
	f(X) = f(X_0) + J_f(X_0)\Delta X + \frac12\Delta X\transp H_f(X_\theta)\Delta X.
\end{align}
\paragraph{带Peano余项的2阶Taylor公式}
设$f$在$X_0$附近属于$\cont^2$，则
\begin{align}
	f(X) = f(X_0) + J_f(X_0)\Delta X + \frac12\Delta X\transp H_f(X_0)\Delta X + \alpha(\Delta X).
\end{align}
当$\Delta X \to 0$时，余项满足
\[
	\alpha(\Delta X) = \frac12\Delta X\transp
	\bigb{H_f(X_0 + \theta\Delta X) - H_f(X_0)}\Delta X = o(\norm{\Delta X}^2).
\]
\paragraph{$m$阶Taylor公式}
为简写高阶导数，记
\begin{equation}\label{eq:calc-directional-taylor-operator}
	D_{\Delta X} := \sum_{i=1}^n\Delta x_i\p_{x_i}.
\end{equation}
求导时将$\Delta X$视为常向量，先对函数求导，再在指定点取值。若$f \in \cont^{m+1}$于包含线段$[X_0,X]$的开邻域，则
\begin{equation}\label{eq:calc-taylor-lagrange}
    f(X) = \sum_{k=0}^m\frac{D_{\Delta X}^kf(X_0)}{k!}
    + \frac{D_{\Delta X}^{m+1}f(X_\theta)}{(m + 1)!},\qquad0 < \theta < 1.
\end{equation}
若只假定$f$在$X_0$附近属于$\cont^m$，则有Peano形式
\begin{equation}\label{eq:calc-taylor-peano}
    f(X) = \sum_{k=0}^m\frac{D_{\Delta X}^kf(X_0)}{k!} + o(\norm{\Delta X}^m).
\end{equation}
以下取$f \in \cont^3$，例式中的偏导均在$(x_0,y_0)$处取值。
$\color{gray}\text{例：}f(x,y) =$
\begin{align*}
	\color{gray}f(x_0,y_0) + \pv fx\Delta x + \pv fy\Delta y + \frac12\biggp{\frac{\p^2f}{\p x^2}\Delta x^2 + 2\pw fxy\Delta x\Delta y + \pv[2]fy\Delta y^2} \\
	\color{gray} + \frac16\biggp{\pv[3]fx\Delta x^3 + 3\frac{\p^3f}{\p x^2\p y}\Delta x^2\Delta y + 3\frac{\p^3f}{\p x\p y^2}\Delta x\Delta y^2 + \pv[3]fy\Delta y^3} + o(\norm{\Delta X}^3).
\end{align*}
$\color{gray}\text{最后换回来：}\Delta x = x - x_0,\Delta y = y - y_0$
\section{极值与条件极值}
\begin{theorem}{Fermat定理}{}
	\begin{center}
		$X_0$是定义域内点，$f$在$X_0$可微，且$X_0$是$f$局部极值点
		\vimplies
		$X_0$是$f$驻点，即$\pv f{x_i}(X_0) = 0,\quad i = 1,2,\ldots,n.$
	\end{center}
\end{theorem}
\begin{theorem}{极值点的充分条件}{}
	$f$在内点$X_0$的邻域内二阶连续可微，且$X_0$是$f$驻点
	\begin{compactenum}[(1)]
		\item $H_f(X_0)$正定$\implies$极小;
		\item $H_f(X_0)$负定$\implies$极大;
		\item $H_f(X_0)$不定$\implies$不是极值点.
	\end{compactenum}
	判断正负定方法：(1)顺序主子式; (2)特征值。半正定或半负定时此判据不能决定是否有极值，还需考察高阶项或沿不同路径的变化。
\end{theorem}
\begin{theorem}{Lagrange乘数法}{}
	设$f,\varphi_1,\ldots,\varphi_{n-k}$均属于$\cont^1$，约束梯度在候选点$X_0$线性无关，从而约束集在该点附近构成$k$维光滑曲面
	\[
		S = \set X{\varphi_i(X) = 0,i = 1,2,\ldots,n - k}.
	\]
	定义Lagrange函数
	\begin{align}
		L(X,\varLambda) := f(X) + \sum_{i=1}^{n-k}\lambda_i\varphi_i(X).
	\end{align}
	若$X_0 \in S$为$f$在$S$上的条件极值点，则存在$\varLambda$使得$(X_0,\varLambda)$为$L$的驻点．
\end{theorem}
\paragraph{计算紧区域或紧曲面上的最值}
连续函数在非空紧集上取得最大值和最小值。对有内部的闭有界区域，应比较内部驻点、各光滑边界部分的约束驻点，以及角点等非光滑点的函数值；对只有曲面的约束集，则直接在曲面上使用约束条件。若约束梯度退化，须另行分析，不能只解Lagrange方程后便排除这些点。
\newpage
